%!TEX root = ../main.tex

\section{Swaps}

\begin{enumerate}
    \item \textbf{Swap:} An OTC agreement between two parties to exchange cashflows in future.
    \item Forward contract is a simple example of swap.
    \item \textbf{Interest Rate Swap:} A swap where interest at a predetermined fixed rate, applied to a certain principal, is exchanged for interest at a floating reference rate, applied to the same principal, with regular exhanges being made for an agreed period of time.
    \item \textbf{Overnight Idexed Swap:} An OIS is an agreement to exchange a fixed rate of interest for a reference rate of interest that is calculated from realized overnight rates.
    \item \textbf{Reasons for trading Interest Rate Swaps:}
    \begin{enumerate}
        \item \textbf{Using the swap to transform a liability}
        \item \textbf{Using the swap to transform an asset}
    \end{enumerate}
    \item \textbf{Comparitive Advantage Argument:} Some companies have a comparitive advantage when borrowing in fixed-rate markets, whereas other companies have a comparitive advantage when borrowing in floating-rate markets. This diffence can be exploited by two firms favoured in different markets by getting into an IRS.
    \item \textbf{Valuation of IRS:} An IRS is worth close to zero when it is first initiated. But it is liable to change as time passes. It can be valued by considering the present value of the future cashflows using forward rates.
    \item \textbf{Fixed-for-Fixed Currency Swaps:} Exchanging principal and interest payments at a fixed rate in one currency for principal and interest payments at a fixed rate in another country.
    \item \textbf{Other Currency Swaps:}
    \begin{enumerate}
        \item \textbf{Fixed-for-floating}
        \item \textbf{Floating-for-floating}
    \end{enumerate}
    \item \textbf{Credit Default Swaps:} A CDS is like an insurance contract that pays off if a particular company or country defaults. The company or the country is known as \textit{reference entity}.
\end{enumerate}